% unit: u04
% sources: sources/080-multiplication-and-inverse-matrices.txt, sources/127-multiplication-and-inverse-matrices.txt, sources/128-mit18-06scf11-ses1-3sum-pdf-download-file.txt, sources/129-inverse-matrices.txt, sources/130-mit18-06scf11-ses1-3prob-pdf-download-file.txt, sources/131-mit18-06scf11-ses1-3sol-pdf-download-file.txt
\section{Multiplication and inverse matrices}\label{sec:u04}

We have been multiplying matrices since Section~\ref{sec:u03}, where elimination steps were
matrices acting on the left. This section sets down the rule for the product $AB$ and then
looks at the same product in four more ways: by columns, by rows, as a sum of
column-times-row pieces, and by blocks. All five give the same answer, and each is the right
one in some argument later on. The second half answers the most important question about a
square matrix: does it have an inverse? A singular matrix sends some nonzero vector to zero,
which no inverse can undo; for an invertible matrix, Gauss--Jordan elimination computes
$A^{-1}$.

\subsection{The rule: row times column}

\begin{definition}[Matrix product]\label{def:u04-product}
Let $A$ be $m\times n$ and $B$ be $n\times p$. The product $C=AB$ is the $m\times p$ matrix
with entries
\[
  c_{ij}=\sum_{k=1}^{n}a_{ik}b_{kj}=a_{i1}b_{1j}+a_{i2}b_{2j}+\cdots+a_{in}b_{nj},
\]
the dot product of row $i$ of $A$ with column $j$ of $B$. The product is defined only when
the number of columns of $A$ equals the number of rows of $B$.
\end{definition}

For instance, the entry $c_{34}$ comes from row~3 of $A$ and column~4 of $B$:
$c_{34}=a_{31}b_{14}+a_{32}b_{24}+\cdots+a_{3n}b_{n4}$. The shapes follow the pattern
$(m\times n)(n\times p)=(m\times p)$: the inner numbers must match and disappear, the outer
numbers give the shape. When $B$ is a single column ($p=1$), this is the row form of
$A\vb{x}$ from Proposition~\ref{prop:u01-rows}.

\subsection{Columns, rows, and columns times rows}

The entry rule hides what the product does to whole vectors. The next three views each look at
whole columns or rows.

\begin{proposition}[Three more views of $AB$]\label{prop:u04-views}
Let $A$ be $m\times n$ with columns $\vb{a}_1,\dots,\vb{a}_n$ and let $B$ be $n\times p$ with
columns $\vb{b}_1,\dots,\vb{b}_p$ and rows $\vb{r}_1\T,\dots,\vb{r}_n\T$.
\begin{enumerate}[(a)]
  \item \emph{Columns:} column $j$ of $AB$ is $A\vb{b}_j$. Every column of $AB$ is a
    combination of the columns of $A$.
  \item \emph{Rows:} row $i$ of $AB$ is (row $i$ of $A$) times $B$. Every row of $AB$ is a
    combination of the rows of $B$.
  \item \emph{Columns times rows:} $AB=\vb{a}_1\vb{r}_1\T+\vb{a}_2\vb{r}_2\T+\cdots+\vb{a}_n\vb{r}_n\T$,
    a sum of $n$ matrices, each an $m\times1$ column times a $1\times p$ row.
\end{enumerate}
\end{proposition}

\begin{proof}
(a) The entries of $A\vb{b}_j$ are $\sum_k a_{ik}b_{kj}$ for $i=1,\dots,m$
(Proposition~\ref{prop:u01-rows}), which is column $j$ of $AB$. By
Definition~\ref{def:u01-Ax}, $A\vb{b}_j=b_{1j}\vb{a}_1+\cdots+b_{nj}\vb{a}_n$.
(b) Entry $j$ of (row $i$ of $A$)$\,B$ is $\sum_k a_{ik}b_{kj}=c_{ij}$, and by
Proposition~\ref{prop:u03-rowtimes} that row is $a_{i1}\vb{r}_1\T+\cdots+a_{in}\vb{r}_n\T$.
(c) The column-times-row product $\vb{a}_k\vb{r}_k\T$ is the $m\times p$ matrix with $(i,j)$
entry $a_{ik}b_{kj}$. Summing over $k$ gives $\sum_k a_{ik}b_{kj}=c_{ij}$.
\end{proof}

\begin{example}[One product, four ways]\label{ex:u04-four-ways}
Let $A=\begin{bmatrix} 1 & 2\\ 3 & 4\end{bmatrix}$ ($2\times2$) and
$B=\begin{bmatrix} 1 & 0 & 2\\ 1 & 1 & 0\end{bmatrix}$ ($2\times3$), so $AB$ is $2\times3$.
\begin{itemize}
  \item \emph{Entries.} $c_{23}=(3,4)\cdot(2,0)=6$; similarly for the other five entries.
  \item \emph{Columns.} $A(1,1)=(1+2,\,3+4)=(3,7)$, $A(0,1)=(2,4)$, $A(2,0)=(2,6)$.
  \item \emph{Rows.} $(1,2)B=1\,(1,0,2)+2\,(1,1,0)=(3,2,2)$ and
    $(3,4)B=3\,(1,0,2)+4\,(1,1,0)=(7,4,6)$.
  \item \emph{Columns times rows.}
    \[
      \begin{bmatrix} 1\\ 3\end{bmatrix}\begin{bmatrix} 1 & 0 & 2\end{bmatrix}
      +\begin{bmatrix} 2\\ 4\end{bmatrix}\begin{bmatrix} 1 & 1 & 0\end{bmatrix}
      =\begin{bmatrix} 1 & 0 & 2\\ 3 & 0 & 6\end{bmatrix}
      +\begin{bmatrix} 2 & 2 & 0\\ 4 & 4 & 0\end{bmatrix}.
    \]
\end{itemize}
Every view gives $AB=\begin{bmatrix} 3 & 2 & 2\\ 7 & 4 & 6\end{bmatrix}$. The
same twelve multiplications are being done each time; only the grouping changes.
\end{example}

A single column times a single row is a very special matrix:
\[
  \begin{bmatrix} 2\\ 3\\ 4\end{bmatrix}\begin{bmatrix} 1 & 6\end{bmatrix}
  =\begin{bmatrix} 2 & 12\\ 3 & 18\\ 4 & 24\end{bmatrix}.
\]
Its columns are multiples of $(2,3,4)$ and its rows are multiples of $(1,6)$, as views (a)
and (b) predict. Plotted as points, the rows $(2,12)$, $(3,18)$, $(4,24)$ all lie on the line
through $(1,6)$, and both columns lie on the line through $(2,3,4)$. In later language
(Section~\ref{sec:u11}), its row space and column space are both lines: it is a matrix of
rank one. View (c) says that every product $AB$ is a sum of such rank-one pieces. For
instance, if $A$ has columns $(2,3,4)$ and $(7,8,9)$ and $B$ has rows $(1,6)$ and $(0,0)$,
the second piece is zero and $AB$ is the matrix above.

\subsection{Block multiplication}

Matrices can be cut into blocks, and if the cuts match, blocks multiply like numbers
(keeping their order). With $A$ and $B$ cut into $2\times2$ arrays of blocks,
\[
  \begin{bmatrix} A_1 & A_2\\ A_3 & A_4\end{bmatrix}
  \begin{bmatrix} B_1 & B_2\\ B_3 & B_4\end{bmatrix}
  =\begin{bmatrix} A_1B_1+A_2B_3 & A_1B_2+A_2B_4\\ A_3B_1+A_4B_3 & A_3B_2+A_4B_4\end{bmatrix},
\]
provided the columns of $A$ are cut the same way as the rows of $B$, so that every product
$A_iB_j$ is defined.

\begin{proposition}[Block multiplication]\label{prop:u04-blocks}
If the column partition of $A$ matches the row partition of $B$, then block $(I,J)$ of $AB$
is $\sum_K A_{IK}B_{KJ}$, a sum of block products.
\end{proposition}

\begin{proof}
Take an entry $c_{ij}$ inside block $(I,J)$. Its sum $\sum_k a_{ik}b_{kj}$ runs over all
column indices $k$ of $A$. Split this sum according to the column blocks $K$: the part with
$k$ in block $K$ is exactly the $(i,j)$ entry of $A_{IK}B_{KJ}$. Adding over $K$ gives the
claim.
\end{proof}

\begin{example}[Blocks behave like numbers]\label{ex:u04-blocks}
Let $C=\begin{bmatrix} 1 & 2\\ 3 & 4\end{bmatrix}$ and let $I$, $0$ be the $2\times2$ identity
and zero matrices. By block multiplication,
\[
  \begin{bmatrix} I & 0\\ C & I\end{bmatrix}\begin{bmatrix} I & 0\\ -C & I\end{bmatrix}
  =\begin{bmatrix} I\cdot I+0\cdot(-C) & I\cdot0+0\cdot I\\ C\cdot I+I\cdot(-C) & C\cdot 0+I\cdot I\end{bmatrix}
  =\begin{bmatrix} I & 0\\ 0 & I\end{bmatrix}.
\]
So the $4\times4$ matrix $\begin{bmatrix} I & 0\\ C & I\end{bmatrix}$, which adds combinations
of rows~1--2 to rows~3--4, is undone by subtracting the same combinations, exactly like the
elimination matrices of Proposition~\ref{prop:u03-Einv}.
\end{example}

Matrix multiplication obeys the associative law (Proposition~\ref{prop:u03-assoc}) and the
distributive laws $A(B+C)=AB+AC$, $(A+B)C=AC+BC$ (Exercise~\ref{ex:u04-distributive}), but
in general $AB\neq BA$.

\subsection{Inverse matrices}

\begin{definition}[Inverse, invertible, singular]\label{def:u04-inverse}
A square $n\times n$ matrix $A$ is \emph{invertible} (or \emph{nonsingular}) if there is an
$n\times n$ matrix $A^{-1}$ with
\[
  A^{-1}A=I\qquad\text{and}\qquad AA^{-1}=I .
\]
A square matrix with no inverse is \emph{singular}.
\end{definition}

An inverse, when it exists, is unique: if $BA=I$ and $AC=I$, then
$B=B(AC)=(BA)C=C$. For square matrices it is enough to find a one-sided inverse
(Corollary~\ref{cor:u04-onesided} below). (A rectangular matrix can have an inverse on one side only, as
Section~\ref{sec:u36} explains.)

Whether $A$ is invertible is the first question to ask about a square matrix. The clearest
way to see that a matrix is \emph{not} invertible is to find a vector it sends to zero.

\begin{theorem}[A nonzero solution of $A\vb{x}=\vb{0}$ rules out an inverse]\label{thm:u04-singular}
If $A$ is square and $A\vb{x}=\vb{0}$ for some $\vb{x}\neq\vb{0}$, then $A$ is singular.
\end{theorem}

\begin{proof}
Suppose $A^{-1}$ existed. Multiplying $A\vb{x}=\vb{0}$ on the left by $A^{-1}$ gives
$\vb{x}=A^{-1}A\vb{x}=A^{-1}\vb{0}=\vb{0}$, contradicting $\vb{x}\neq\vb{0}$.
\end{proof}

\begin{example}[A singular matrix, three ways]\label{ex:u04-singular}
Take $A=\begin{bmatrix} 1 & 3\\ 2 & 6\end{bmatrix}$.
\emph{Null vector:} $\vb{x}=(3,-1)$ gives $A\vb{x}=3(1,2)-1(3,6)=(0,0)$, so $A$ is singular
by Theorem~\ref{thm:u04-singular}.
\emph{Columns:} column~2 is $3$ times column~1, so both columns lie on the line through
$(1,2)$. By Proposition~\ref{prop:u04-views}(a), every column of any product $AB$ lies on that
line too, so $AB$ can never have the column $(1,0)$ of $I$.
\emph{Elimination:} subtracting $2$ times row~1 from row~2 leaves $(0,0)$, so there is no
second pivot. (For readers who know determinants: $1\cdot6-3\cdot2=0$.)
\end{example}

\begin{figure}[ht]
\centering
\begin{tikzpicture}[scale=0.55,>=Stealth]
  \draw[->] (-1,0) -- (4.2,0);
  \draw[->] (0,-0.8) -- (0,7.4);
  \draw[dashed,gray] (-0.4,-0.8) -- (3.7,7.4);
  \draw[->,very thick,orange!85!black] (0,0) -- (3,6) node[right] {$(3,6)$};
  \draw[->,very thick,blue] (0,0) -- (1,2) node[right] {$(1,2)$};
  \fill[red!75!black] (1,0) circle (3pt) node[below right] {$(1,0)$};
  \node[font=\small,align=center] at (1.6,-2.6) {$\begin{bmatrix} 1&3\\ 2&6\end{bmatrix}$: columns on one line,\\ $(1,0)$ out of reach};
\end{tikzpicture}
\qquad\qquad
\begin{tikzpicture}[scale=0.3,>=Stealth]
  \draw[->] (-1.5,0) -- (8.5,0);
  \draw[->] (0,-1.5) -- (0,15);
  \draw[->,thick,dashed,blue] (0,0) -- (7,14) node[right] {$7(1,2)$};
  \draw[->,thick,dashed,orange!85!black] (7,14) -- (1,0);
  \node[orange!85!black,right] at (4.3,6.5) {$-2(3,7)$};
  \draw[->,very thick,orange!85!black] (0,0) -- (3,7);
  \draw[->,very thick,blue] (0,0) -- (1,2);
  \fill[red!75!black] (1,0) circle (6pt) node[below right] {$(1,0)$};
  \node[font=\small,align=center] at (3.5,-4.2) {$\begin{bmatrix} 1&3\\ 2&7\end{bmatrix}$: independent columns,\\ $7(1,2)-2(3,7)=(1,0)$};
\end{tikzpicture}
\caption{Left: the columns of the singular matrix lie on one line (dashed), so every
combination stays on it and the column $(1,0)$ of $I$ cannot be produced. Right: the columns
$(1,2)$ and $(3,7)$ of the invertible matrix (solid, near the origin) point in different
directions, so their combinations fill the plane. Seven copies of the first column followed
by $-2$ copies of the second (dashed) land on $(1,0)$; the weights $7,-2$ form the first column
of $A^{-1}$. They are large because the two columns are nearly parallel.}
\label{fig:u04-columns}
\end{figure}

\subsection{Gauss--Jordan elimination computes \texorpdfstring{\(A^{-1}\)}{the inverse}}

Now change the $6$ to a $7$: $A=\begin{bmatrix} 1 & 3\\ 2 & 7\end{bmatrix}$. Its columns point
in different directions (Figure~\ref{fig:u04-columns}), and we look for
$A^{-1}=\begin{bmatrix} a & c\\ b & d\end{bmatrix}$ with $AA^{-1}=I$. By the column view,
this is two systems with the same matrix:
\[
  A\begin{bmatrix} a\\ b\end{bmatrix}=\begin{bmatrix} 1\\ 0\end{bmatrix},\qquad
  A\begin{bmatrix} c\\ d\end{bmatrix}=\begin{bmatrix} 0\\ 1\end{bmatrix}.
\]
In general, column $j$ of $A^{-1}$ solves $A\vb{x}=\vb{e}_j$, where $\vb{e}_j$ is column $j$
of $I$. Finding $A^{-1}$ means solving $n$ systems at once. Gauss's elimination would solve
them one at a time; Jordan's idea is to carry all $n$ right-hand sides along as extra columns
and keep eliminating \emph{upwards} until the left half is $I$.

\begin{example}[Gauss--Jordan on a $2\times2$ matrix]\label{ex:u04-gj2}
\[
  \left[\begin{array}{rr|rr} 1 & 3 & 1 & 0\\ 2 & 7 & 0 & 1\end{array}\right]
  \xrightarrow{\ \text{row}_2-2\,\text{row}_1\ }
  \left[\begin{array}{rr|rr} 1 & 3 & 1 & 0\\ 0 & 1 & -2 & 1\end{array}\right]
  \xrightarrow{\ \text{row}_1-3\,\text{row}_2\ }
  \left[\begin{array}{rr|rr} 1 & 0 & 7 & -3\\ 0 & 1 & -2 & 1\end{array}\right].
\]
The first step is ordinary elimination (Gauss stops at the upper triangular form). The second
clears the entry \emph{above} the pivot (Jordan continues). The right half is
$A^{-1}=\begin{bmatrix} 7 & -3\\ -2 & 1\end{bmatrix}$. Check:
$A^{-1}A=\begin{bmatrix} 7-6 & 21-21\\ -2+2 & -6+7\end{bmatrix}=I$.
\end{example}

\begin{theorem}[Why Gauss--Jordan works]\label{thm:u04-gj}
Suppose a sequence of row operations (elimination steps, row exchanges, and division of rows
by nonzero pivots) turns $[\,A\;\;I\,]$ into $[\,I\;\;E\,]$. Then $E=A^{-1}$: both $EA=I$
and $AE=I$.
\end{theorem}

\begin{proof}
Each row operation is multiplication on the left by an invertible matrix (an elimination
matrix, a permutation, or a diagonal matrix with one nonzero entry changed), so the whole
sequence is multiplication by their product, call it $M$. By block multiplication
$M[\,A\;\;I\,]=[\,MA\;\;M\,]$, and this equals $[\,I\;\;E\,]$. So $M=E$ and $EA=I$.
Moreover $M$ is invertible, its inverse being the product of the inverse steps in reverse
order. Hence $A=M^{-1}(MA)=M^{-1}$, and $AE=M^{-1}M=I$.
\end{proof}

\begin{example}[A $3\times3$ inverse]\label{ex:u04-gj3}
Invert $K=\begin{bmatrix} 2 & -1 & 0\\ -1 & 2 & -1\\ 0 & -1 & 2\end{bmatrix}$. Forward
elimination (multipliers $-\tfrac12$, then $-\tfrac23$):
{\renewcommand{\arraystretch}{1.3}
\[
  \left[\begin{array}{rrr|rrr} 2&-1&0&1&0&0\\ -1&2&-1&0&1&0\\ 0&-1&2&0&0&1\end{array}\right]
  \to
  \left[\begin{array}{rrr|rrr} 2&-1&0&1&0&0\\ 0&\tfrac32&-1&\tfrac12&1&0\\ 0&0&\tfrac43&\tfrac13&\tfrac23&1\end{array}\right].
\]}%
Upward elimination: add $\tfrac34$ of row~3 to row~2, then $\tfrac23$ of the new row~2 to
row~1; finally divide by the pivots $2,\tfrac32,\tfrac43$:
{\renewcommand{\arraystretch}{1.3}
\[
  \to
  \left[\begin{array}{rrr|rrr} 2&0&0&\tfrac32&1&\tfrac12\\ 0&\tfrac32&0&\tfrac34&\tfrac32&\tfrac34\\ 0&0&\tfrac43&\tfrac13&\tfrac23&1\end{array}\right]
  \to
  \left[\begin{array}{rrr|rrr} 1&0&0&\tfrac34&\tfrac12&\tfrac14\\ 0&1&0&\tfrac12&1&\tfrac12\\ 0&0&1&\tfrac14&\tfrac12&\tfrac34\end{array}\right].
\]}%
So $K^{-1}=\frac14\begin{bmatrix} 3&2&1\\ 2&4&2\\ 1&2&3\end{bmatrix}$. Notice that $K$ is
symmetric and so is $K^{-1}$, and that $K^{-1}$ has no zero entries although $K$ has two.
\end{example}

Putting the pieces together gives the basic test for invertibility of a square matrix.

\begin{theorem}[Invertibility test]\label{thm:u04-test}
For an $n\times n$ matrix $A$ the following are equivalent:
\begin{enumerate}[(a)]
  \item $A$ is invertible;
  \item $A\vb{x}=\vb{0}$ has only the solution $\vb{x}=\vb{0}$;
  \item elimination, with row exchanges allowed, produces $n$ nonzero pivots.
\end{enumerate}
\end{theorem}

\begin{proof}
(a)$\Rightarrow$(b) is Theorem~\ref{thm:u04-singular}. (c)$\Rightarrow$(a): with $n$ pivots,
continue upward and divide by the pivots; this is Gauss--Jordan, and
Theorem~\ref{thm:u04-gj} produces $A^{-1}$. (b)$\Rightarrow$(c): suppose elimination stops at
column $k$ because the entries of the current matrix $M$ in column $k$, from row $k$ down, are
all zero. The first $k-1$ columns of $M$ have pivots in rows $1,\dots,k-1$ and zeros below,
so column $k$ of $M$ is a combination of them: back substitution in the top
$(k-1)\times(k-1)$ triangle (Proposition~\ref{prop:u03-back}) gives weights
$y_1,\dots,y_{k-1}$. Then $\vb{x}=(-y_1,\dots,-y_{k-1},1,0,\dots,0)\neq\vb{0}$ satisfies
$M\vb{x}=\vb{0}$. Since $M$ is $A$ multiplied on the left by invertible step matrices, also
$A\vb{x}=\vb{0}$, contradicting (b).
\end{proof}

\begin{corollary}[One side is enough]\label{cor:u04-onesided}
If $A$ and $B$ are $n\times n$ and $BA=I$, then $A$ is invertible, $B=A^{-1}$, and $AB=I$.
\end{corollary}

\begin{proof}
If $A\vb{x}=\vb{0}$ then $\vb{x}=BA\vb{x}=B\vb{0}=\vb{0}$, so (b) of
Theorem~\ref{thm:u04-test} holds and $A^{-1}$ exists. Then
$B=B(AA^{-1})=(BA)A^{-1}=A^{-1}$, and $AB=AA^{-1}=I$. (If instead $AB=I$, apply this to the
pair $B,A$.)
\end{proof}

\begin{trap}
The entry $c_{ij}$ of $AB$ is \emph{not} $a_{ij}b_{ij}$; it uses a whole row of $A$ and a whole
column of $B$. And the right half of $[\,A\;\;I\,]$ becomes $A^{-1}$ only if the left half
really reaches $I$; stopping at the triangular form $U$ leaves a matrix that is not the
inverse. Finally, a matrix with no zero entries can still be singular, as
$\begin{bmatrix} 1&3\\ 2&6\end{bmatrix}$ shows.
\end{trap}

\subsection{Exercises}

\begin{exercise}\label{ex:u04-distributive}
With $A=\begin{bmatrix} 1&2\\ 3&4\end{bmatrix}$, $B=\begin{bmatrix} 1&0\\ 0&0\end{bmatrix}$
and $C=\begin{bmatrix} 0&0\\ 5&6\end{bmatrix}$, compute $AB+AC$ and $A(B+C)$. Then prove
$A(B+C)=AB+AC$ in general.
\end{exercise}
\begin{solution}
$AB=\begin{bmatrix} 1&0\\ 3&0\end{bmatrix}$ (column~1 of $A$ in column~1) and
$AC=\begin{bmatrix} 10&12\\ 20&24\end{bmatrix}$ (rows of $AC$ are $2(5,6)$ and $4(5,6)$), so
$AB+AC=\begin{bmatrix} 11&12\\ 23&24\end{bmatrix}$. Also
$B+C=\begin{bmatrix} 1&0\\ 5&6\end{bmatrix}$ and
$A(B+C)=\begin{bmatrix} 1+10&0+12\\ 3+20&0+24\end{bmatrix}=\begin{bmatrix} 11&12\\ 23&24\end{bmatrix}$,
the same. In general, the $(i,j)$ entry of $A(B+C)$ is
$\sum_k a_{ik}(b_{kj}+c_{kj})=\sum_k a_{ik}b_{kj}+\sum_k a_{ik}c_{kj}$, which is the $(i,j)$
entry of $AB+AC$.
\end{solution}

\begin{exercise}\label{ex:u04-Uinv}
Use Gauss--Jordan elimination on $[\,U\;\;I\,]$ to find the inverse of the upper triangular
matrix $U=\begin{bmatrix} 1&a&b\\ 0&1&c\\ 0&0&1\end{bmatrix}$.
\end{exercise}
\begin{solution}
$U$ is already triangular with pivots $1,1,1$, so only upward steps are needed. Subtract $a$
times row~2 from row~1 and $c$ times row~3 from row~2:
\[
  \left[\begin{array}{ccc|ccc} 1&a&b&1&0&0\\ 0&1&c&0&1&0\\ 0&0&1&0&0&1\end{array}\right]
  \to
  \left[\begin{array}{ccc|ccc} 1&0&b-ac&1&-a&0\\ 0&1&0&0&1&-c\\ 0&0&1&0&0&1\end{array}\right].
\]
Then subtract $b-ac$ times row~3 from row~1:
\[
  \to
  \left[\begin{array}{ccc|ccc} 1&0&0&1&-a&ac-b\\ 0&1&0&0&1&-c\\ 0&0&1&0&0&1\end{array}\right],
  \qquad
  U^{-1}=\begin{bmatrix} 1&-a&ac-b\\ 0&1&-c\\ 0&0&1\end{bmatrix}.
\]
Check the $(1,3)$ entry of $UU^{-1}$: $1\cdot(ac-b)+a\cdot(-c)+b\cdot1=0$. The inverse of an
upper triangular matrix is upper triangular.
\end{solution}

\begin{exercise}\label{ex:u04-ab}
For which numbers $a$ and $b$ is $A=\begin{bmatrix} a&b&b\\ a&a&b\\ a&a&a\end{bmatrix}$
invertible? Find $A^{-1}$ when it exists.
\end{exercise}
\begin{solution}
Quick tests first: if $a=0$ the first column is zero, and if $a=b$ all rows are equal; either
way $A$ is singular. Gauss--Jordan: subtract row~1 from rows~2 and~3, then the new row~2 from
row~3:
\[
  \left[\begin{array}{ccc|ccc} a&b&b&1&0&0\\ a&a&b&0&1&0\\ a&a&a&0&0&1\end{array}\right]
  \to
  \left[\begin{array}{ccc|ccc} a&b&b&1&0&0\\ 0&a-b&0&-1&1&0\\ 0&0&a-b&0&-1&1\end{array}\right].
\]
The pivots are $a$, $a-b$, $a-b$, so $A$ is invertible exactly when $a\neq0$ and $a\neq b$.
Then divide rows~2 and~3 by $a-b$, divide row~1 by $a$, and subtract $\tfrac ba$ times rows~2
and~3 from row~1. The right half becomes
\[
  A^{-1}=\frac{1}{a-b}\begin{bmatrix} 1&0&-b/a\\ -1&1&0\\ 0&-1&1\end{bmatrix}.
\]
For example, the $(1,1)$ entry is $\frac1a+\frac ba\cdot\frac{1}{a-b}=\frac{(a-b)+b}{a(a-b)}=\frac{1}{a-b}$.
Check row~3 of $A$ times column~1 of $A^{-1}$: $\frac{a\cdot1+a\cdot(-1)+a\cdot0}{a-b}=0$.
\end{solution}

\begin{exercise}\label{ex:u04-rank1}
Write $\begin{bmatrix} 1&2\\ 3&4\end{bmatrix}\begin{bmatrix} 5&6\\ 7&8\end{bmatrix}$ as a sum of
two column-times-row matrices, and check the result against the entry rule.
\end{exercise}
\begin{solution}
By Proposition~\ref{prop:u04-views}(c),
\[
  \begin{bmatrix} 1\\ 3\end{bmatrix}\begin{bmatrix} 5&6\end{bmatrix}
  +\begin{bmatrix} 2\\ 4\end{bmatrix}\begin{bmatrix} 7&8\end{bmatrix}
  =\begin{bmatrix} 5&6\\ 15&18\end{bmatrix}+\begin{bmatrix} 14&16\\ 28&32\end{bmatrix}
  =\begin{bmatrix} 19&22\\ 43&50\end{bmatrix}.
\]
Entry rule: $c_{11}=1\cdot5+2\cdot7=19$, $c_{12}=1\cdot6+2\cdot8=22$,
$c_{21}=3\cdot5+4\cdot7=43$, $c_{22}=3\cdot6+4\cdot8=50$.
\end{solution}
